% 完整英文前言 intro_en.tex 的中文初译；待同底本精修。
\begin{abstract}
若平面星形集在每个无向方向上都包含一条闭单位线段，其 Lebesgue 外测度至少为 $1/10$。本文给出完整证明，不要求原集合可测、有界，也不假设能够规则地选择其中的针。证明测量实际的圆周截面，仅在近端点半径允许时使用第二个端点，并在同一份实际面积预算上平衡其余方向。随后，我们研究完整单位针的上界构造及其严格改进，再考察极值问题中若干重构方法的障碍。我们构造了一个固定的有限竞争星体 $K$：对每个同中心且容纳随方向连续、首尾闭合的单位针选择的紧星体 $Q$，都有 $|Q\setminus K|\ge\varepsilon_K>0$，而 $Q$ 的外半径可以任意大。一个规定的双剪切规则，也会在趋近面积下确界的有限输入上严格增加面积。有限采样、半径截断、最长完整弦及强面积闭合，进一步说明这些方法保留了什么信息。锐下确界仍未确定；缺少的比较必须把完整针的可行性与高效构造中的重叠关系结合起来。
\end{abstract}

\section{同一个中心与独立放置的针}\label{sec:introduction}
平面 Kakeya 条件要求每个方向上都有一条单位线段。不加其他几何要求时，这与任意小的面积相容。星形性增加了一项要求：必须存在一个点，使它与集合中每一点之间的线段都仍在该集合中。因此，每条可用的针都会带来径向材料。问题在于，独立放置这些针时，材料究竟能够重叠到什么程度。

Cunningham 在对星形 Kakeya 集的研究中证明了正的面积下界 $\pi/108$~\cite{Cunningham1971}。他在第 129 页的文末说明中指出，这个下界证明只使用每个方向存在一条单位线段，并未使用连续反转过程。Li 随后通过圆周截面估计，将这一静态问题的外测度下界提高到 $\pi/98$~\cite[定理~1.1]{Li2026}。本文沿用对实际圆周截面进行估计的思路，在内圆盘中进一步结合近端点资格与互补权重。连续闭合的针选择会另行明确引入，而不会作为任意 Kakeya 集的默认条件。

我们的第一个结果是
\[
 E\subset\mathbb R^2\text{ 为星形集且每个方向都有完整单位针}
 \quad\Longrightarrow\quad |E|^*\ge\frac1{10}.
\]
这里 $|E|^*$ 是 Lebesgue 外测度，星形中心可以是任意点。原针的选择可以完全不规则。转到任意开超集后，可以在该开集中重新选出可测族，再对实际的极坐标截面积分。所用的可测投影事实属于标准描述测度论~\cite{Kechris1995}；几何估计及完整的面积分配都在下文证明。

证明的低半径部分有一项具体的补偿关系。近端点离中心足够远的方向，提供两个不同的理想角度锚点；缺少第二个锚点的方向，则有半径更大的远端点，能在稍后的径向区间上给出更强的下界。对这两条估计作凸组合，可以让每个实际点至多贡献一次，并使最终的下界系数不依赖于两类方向各占多少。其余更大的端点尺度由互不相交的外侧径向区域处理。整个过程保留了零高度、垂足在底边外的情况、闭端点，以及实际角度与无向方向的区别。

上界构造从另一侧使用相同的几何对象。它们安排完整底边，使一个端点新增的材料被已有材料遮住，再利用仍然暴露的部分证明严格改进。这些改进必须在微观极限和有限记录恢复两步之后仍保持为正。因此，最后的伸长操作在原微观尺度下控制法向位移，并检查整条内接弦。它的比较适用于指定的源族，而不直接适用于任意近极小输入。

论文的其余部分考察哪些方法可以把这些构造传递到全体输入。两种方法因不同的几何原因失效：几乎保留一个固定星体的全部材料，可能与连续针选择不相容；固定剪切规则的闭轴分支，则可以积累出真正的新增面积。对空间逃逸和极限代表，则能给出正面结论。远处的针允许带有定量误差的半径截断；强面积度量允许构造实际的全方向代表，并给出全局相对替换不等式。这两项结果都尚未提供近极小序列的紧性或所缺少的锐值比较。区分它们的作用，才能明确后续论证需要增加什么。
